\subsection{
    Taxonomia de Vulnerabilidades
} \label{sec:taxonomia_vulnerabilidades}


A taxonomia de vulnerabilidades permite organizar falhas de segurança em categorias estruturadas, facilitando tanto a análise quanto a mitigação desses problemas. Nesse contexto, diferentes \textit{frameworks} e organizações especializadas têm desenvolvido modelos de classificação amplamente utilizados pela comunidade de segurança da informação. Entre eles, destaca-se o trabalho da \textit{OWASP Foundation}, que mantém uma das referências mais utilizadas para identificação e priorização de vulnerabilidades em aplicações web. 

Mecanismos de segurança correspondem a técnicas e controles para prevenir, detectar e mitigar falhas, como validação de entrada, autenticação multifator, criptografia, controle de acesso e monitoramento de \textit{logs}. Arquiteturas de segurança, por sua vez, descrevem como esses mecanismos são integrados no sistema, estabelecendo camadas de proteção. Estratégias como \textit{defense in depth} e isolamento de componentes em microsserviços exemplificam esse arranjo em aplicações modernas.

\subsubsection{
    \textit{Framework OWASP Top 10}: evolução 2017 → 2021
}

Entre 2017 e 2021, o ranking do OWASP Top 10 foi reordenado. \textit{Broken Access Control} passou ao primeiro lugar, evidenciando a necessidade de controles de autorização mais rigorosos em todo o ciclo de desenvolvimento. Falhas criptográficas subiram para a segunda posição, enquanto injeção caiu para a terceira (incidência máxima de 19\% e média de 3,37\%). Além disso, a categoria \textit{Insecure Design} reforçou a importância de decisões arquiteturais seguras desde fases iniciais, evitando dependência de validações apenas no \textit{frontend} \cite{OWASP2021Top10}.

\subsubsection{
    Vulnerabilidades de Injeção (\textit{SQLi, Command, LDAP})
}

Vulnerabilidades de injeção ocorrem quando entradas do usuário são interpretadas como comandos por subsistemas, como banco de dados, sistema operacional ou diretórios LDAP. A mitigação requer validação e sanitização rigorosas, uso de abstrações seguras (como ORM), princípio do menor privilégio e monitoramento contínuo de \textit{logs}. Em \textit{SQL injection (SQLi)}, entradas manipuladas alteram consultas (por exemplo, com tautologias como \verb|OR 1=1|), permitindo leitura ou alteração indevida de dados \cite{Appelt2014SQLInjectionTesting}. Já \textit{command injection} pode levar à execução arbitrária no sistema operacional, e \textit{LDAP injection} permite alterar a lógica de consultas a diretórios, resultando em acesso não autorizado.

\subsubsection{
    \textit{Cross-Site Scripting (XSS): Reflected, Stored, DOM-based}
}
\textit{Cross-Site Scripting (XSS)} permite a injeção de \textit{scripts} maliciosos em páginas visualizadas por outros usuários, quando dados de entrada não são devidamente tratados \cite{Sarmah2018SurveyXSSDetection}. No \textit{Reflected XSS}, o conteúdo malicioso é retornado imediatamente na resposta; no \textit{Stored XSS}, ele é persistido no servidor e executado para múltiplos usuários; no \textit{DOM-based XSS}, a exploração ocorre no navegador, por manipulação insegura do DOM em \textit{JavaScript}. Os impactos incluem roubo de sessão, redirecionamentos maliciosos, execução de ações em nome da vítima e disseminação de \textit{malware}.

\subsubsection{
    \textit{Cross-Site Request Forgery (CSRF)}
}

\textit{Cross-Site Request Forgery (CSRF)} explora a confiança da aplicação no navegador autenticado da vítima. O atacante induz o usuário a disparar uma requisição maliciosa para um sistema onde já possui sessão ativa; como \textit{cookies} de autenticação são enviados automaticamente, o servidor tende a interpretar a ação como legítima \cite{Lin2009ThreatModelingCSRF}. 

\subsubsection{
    \textit{Broken Access Control} e IDOR
}

\textit{Broken Access Control} ocorre quando regras de autorização são aplicadas de forma incompleta, permitindo acesso ou ações além do perfil do usuário \cite{Kusnana2025IDORBrokenAccess}. Por exemplo, um \textit{endpoint} administrativo como \verb|https://site.com/admin/financas| torna-se acessível sem validação adequada de permissões. O \textit{Insecure Direct Object Reference} (IDOR) é um caso recorrente: ao expor identificadores em URLs ou parâmetros sem checagem de autorização, a aplicação permite manipulação de IDs e acesso indevido a recursos de terceiros, como em \verb|https://site.com/pedido?id=1025|.


\subsubsection{
    Falhas Criptográficas
}
Falhas criptográficas decorrem do uso inadequado de mecanismos de proteção de dados sensíveis. Entre as causas mais comuns estão algoritmos obsoletos (como SHA-1, DES e 3DES), gestão incorreta de chaves, geração previsível de números aleatórios e reutilização de vetores de inicialização, fatores que podem comprometer completamente a confidencialidade \cite{Blessing2021RollYourOwnCrypto}. Também são recorrentes o armazenamento inseguro de senhas (\textit{plain text} ou hash inadequado) e erros de transmissão, como ausência de HTTPS ou configuração incorreta de certificados, facilitando ataques \textit{man-in-the-middle} (MitM).


\subsubsection{\textit{Misconfiguration} e \textit{Supply Chain}}

Más configurações (\textit{misconfiguration}) surgem quando sistemas são implantados com parâmetros inseguros, por ausência de padronização ou erro operacional. Exemplos incluem \textit{buckets} públicos, portas administrativas expostas, bancos sem autenticação adequada e credenciais padrão. Em ambientes \textit{cloud}, esse problema permanece entre as principais causas de violação de dados \cite{CloudSecurityAlliance2024TopThreats}. Já ataques de \textit{supply chain} exploram relações de confiança em bibliotecas, \textit{frameworks} e pipelines de CI/CD: ao comprometer um componente amplamente usado, o atacante pode propagar código malicioso em larga escala, como no caso \textit{event-stream} (NPM, 2018), com impacto sobre confidencialidade, integridade e disponibilidade.
